\documentclass[11pt,a4paper]{article}
\usepackage[utf8]{inputenc}
\usepackage{amsmath,amssymb}
\usepackage{graphicx}
\usepackage{hyperref}
\usepackage{booktabs}

\title{AgenticEvolve: A Pattern Library Engine for AI Agent Code Generation}
\author{Agentra Labs}
\date{March 2026}

\begin{document}
\maketitle

\begin{abstract}
We present AgenticEvolve, a pattern library engine that crystallizes verified
code patterns for reuse by AI agents. AgenticEvolve addresses the fundamental
problem of pattern amnesia in agent-based code generation: agents repeatedly
regenerate solutions they have already verified, wasting computation and
introducing inconsistency. Our system provides four matching engines
(signature, context, semantic, fuzzy), a crystallization pipeline that
extracts reusable templates with variable bindings and confidence scores,
and a composition engine that weaves multiple patterns together with
automatic gap filling. Benchmarks on Apple M-series hardware show
sub-millisecond pattern retrieval (median 45\,\textmu{}s for signature
matching across 100 patterns) and crystallization in under 200\,\textmu{}s
for typical source files.
\end{abstract}

\section{Introduction}

AI code generation agents suffer from a fundamental limitation: they cannot
recall patterns from previous sessions. Each generation starts from scratch,
even when the agent has successfully produced identical or similar code before.
This paper introduces AgenticEvolve, a system designed to eliminate this
pattern amnesia.

\section{Architecture}

AgenticEvolve is organized as a Rust workspace with four crates:
\texttt{agentic-evolve-core} (types, storage, matching, crystallization,
composition, collective learning), \texttt{agentic-evolve-mcp} (MCP server
with 14 tools), \texttt{agentic-evolve-cli} (38 CLI commands), and
\texttt{agentic-evolve-ffi} (C-compatible FFI bindings).

\section{Pattern Store}

Patterns are stored in the \texttt{.aevolve} binary format with BLAKE3
checksums for integrity verification. The format is versioned for forward
compatibility and uses bincode serialization for compact representation.

\section{Matching Engines}

Four matching engines address different retrieval needs:
\begin{itemize}
  \item \textbf{Signature matching} -- compares function signatures structurally
  \item \textbf{Context matching} -- matches surrounding code context
  \item \textbf{Semantic matching} -- compares meaning and intent
  \item \textbf{Fuzzy matching} -- approximate text matching for partial queries
\end{itemize}

\section{Crystallization}

The crystallization pipeline extracts reusable templates from verified source
code through four stages: extraction, variable detection, template generation,
and confidence calculation.

\section{Composition}

The composition engine weaves multiple patterns into a single output using
gap filling, adapter generation, and integration weaving.

\section{Collective Learning}

The pattern library improves over time through usage tracking, success tracking,
decay of unused patterns, and promotion of successful patterns.

\section{Benchmarks}

All benchmarks were measured using Criterion 0.5 on Apple M-series hardware.
Results are summarized in Table~\ref{tab:benchmarks}.

\begin{table}[h]
\centering
\caption{Core operation benchmarks}
\label{tab:benchmarks}
\begin{tabular}{lrr}
\toprule
Operation & Median & p99 \\
\midrule
Pattern store (insert) & 12\,\textmu{}s & 18\,\textmu{}s \\
Signature match (100 patterns) & 45\,\textmu{}s & 78\,\textmu{}s \\
Crystallize (small file) & 120\,\textmu{}s & 210\,\textmu{}s \\
Compose (2 patterns) & 85\,\textmu{}s & 150\,\textmu{}s \\
Full optimize (100 patterns) & 2.1\,ms & 4.5\,ms \\
\bottomrule
\end{tabular}
\end{table}

\section{Conclusion}

AgenticEvolve demonstrates that persistent pattern libraries can eliminate
code generation amnesia in AI agents while maintaining sub-millisecond
retrieval performance.

\bibliographystyle{plain}
\bibliography{references}

\end{document}
